\documentclass[11pt]{amsart}
\usepackage[margin=0.92in]{geometry}
\usepackage{amsmath,amssymb,amsthm,mathtools}
\linespread{0.99}
\setlength{\parskip}{2pt}

\newtheorem{theorem}{Theorem}[section]
\newtheorem{proposition}[theorem]{Proposition}
\newtheorem{lemma}[theorem]{Lemma}
\theoremstyle{definition}
\newtheorem{remark}[theorem]{Remark}
\theoremstyle{remark}
\newtheorem*{checkpoint}{\textbf{Open point}}

\newcommand{\Om}{\Omega}
\newcommand{\dT}{\delta_T}
\newcommand{\Asym}{\mathcal A}
\newcommand{\R}{\mathcal R}
\newcommand{\Sg}{\Sigma}
\newcommand{\fbar}{\bar f}
\newcommand{\Hn}{\mathcal H^{n-1}}
\newcommand{\Rn}{\mathbb R^n}
\newcommand{\TF}{\bigl(D^2v+\tfrac1n I\bigr)}
\newcommand{\norm}[1]{\left\lVert#1\right\rVert}
\DeclareMathOperator{\tr}{tr}
\DeclareMathOperator{\dist}{dist}

\title{Route~B, clarified against the near-Serrin picture}
\author{Working note (Plan 3)}
\date{\today}

\begin{document}
\maketitle

\begin{abstract}
You proposed, for the graph route (Strategy~A), to take the small
problematic part of a near-Serrin boundary and \emph{smooth it out} by
passing to an interior level set, using the Poisson kernel (interior
elliptic regularity), and then to bootstrap higher derivatives from the
fact that $|\nabla u|$ varies little. This note explains, in full, what
Route~B does with the \emph{same} input ($|\nabla u|$ close to constant
on a level, in $L^2$), why it is the \emph{opposite} design choice, and
exactly where each route's single hard point sits. Everything is
written out: the level-set defect, the Markov selection, the
Rellich--Pohozaev and Weinberger identities, the trace-free Hessian
rigidity, the small-gradient excision, the Serrin dictionary, and a
worked ball-plus-tentacle illustration.
\end{abstract}

\section{Setup and the common input}

Let $\Om\subset\Rn$ be open, $|\Om|=\omega_n$, with torsion function
$-\Delta u=1$ in $\Om$, $u=0$ on $\partial\Om$, $u>0$ inside. Write
$T(\Om)=\int_\Om|\nabla u|^2=\int_\Om u$, let $B$ be the volume-matched
ball, and set the scale-invariant Saint--Venant deficit and the
Fraenkel asymmetry
\[
  \dT(\Om)=\frac{T(B)-T(\Om)}{T(B)}\ge0,\qquad
  \Asym(\Om)=\min_{x_0}\frac{|\Om\triangle B(x_0)|}{|\Om|}.
\]
The target is the sharp inequality $\Asym(\Om)^2\le C(n)\,\dT(\Om)$.

\subsection{The level-set defect}
For a regular value $t$ write $E_t=\{u>t\}$, $\Sg_t=\{u=t\}$,
$m(t)=|E_t|$, $P(t)=\Hn(\Sg_t)$, and the mean
$\fbar_t=m(t)/P(t)$. From $-\Delta u=1$ one has the two elementary
identities
\[
  m(t)=\int_{E_t}(-\Delta u)=\int_{\Sg_t}|\nabla u|\,d\Hn,
  \qquad
  -m'(t)=\int_{\Sg_t}\frac{d\Hn}{|\nabla u|}.
\]
Cauchy--Schwarz, $P(t)^2\le\bigl(\int_{\Sg_t}|\nabla u|^{-1}\bigr)
\bigl(\int_{\Sg_t}|\nabla u|\bigr)$, defines the
\emph{boundary-gradient defect}
\[
  D_H(t):=\Bigl(\int_{\Sg_t}\tfrac{d\Hn}{|\nabla u|}\Bigr)
          \Bigl(\int_{\Sg_t}|\nabla u|\,d\Hn\Bigr)-P(t)^2\ \ge0,
\]
which vanishes exactly when $|\nabla u|$ is constant on $\Sg_t$ (the
Serrin overdetermined condition). Expanding the square gives the exact
\emph{weighted-variance identity}
\begin{equation}\label{eq:wv}
  \int_{\Sg_t}\frac{(|\nabla u|-\fbar_t)^2}{|\nabla u|}\,d\Hn
  \;=\;\frac{m(t)}{P(t)^2}\,D_H(t).
\end{equation}
Thus $D_H(t)$ is precisely the variance of $|\nabla u|$ on $\Sg_t$
\emph{weighted by} $1/|\nabla u|$; this weight is the protagonist
throughout. The level-set deficit identity (Plan~2) reads
$\dT(\Om)\simeq_n\int (D_H(t)+D_I(t))\,w(t)\,dt$ with a positive weight
$w$ and the isoperimetric defect $D_I\ge0$; only $D_H$ is used here.

\subsection{The Markov selection}
\begin{lemma}\label{lem:markov}
Fix a small absolute constant $\theta\in(0,1)$. There is a level
$\hat t$ with $D_H(\hat t)\le n^2\omega_n\,\theta$ (equivalently, the
left side of \eqref{eq:wv} at $\hat t$ is $\le\theta$), with the
discarded layer $|\Om\setminus E_{\hat t}|=O_\theta(\dT(\Om))$, and
with $\hat t$ a regular value.
\end{lemma}
\begin{proof}[Proof sketch]
Drop $D_I\ge0$ in the deficit identity and change variable to the
volume parameter $s=m(t)$, in which the layer is exactly linear,
$|\Om\setminus E_t|=\omega_n-s$. The deficit bound becomes
$\int_0^{\omega_n}\Phi(s)\,ds\le C_n\dT$ for the transported integrand
$\Phi$; Markov gives $|\{\Phi>\theta'\}|\le C_n\dT/\theta'$, so a good
$\hat t$ exists inside a near-boundary slab of $s$-length
$O_\theta(\dT)$; intersect with the full-measure set of regular values
(Sard). The weight $w$ is two-sidedly $O(1)$ on that slab because
$m\to\omega_n$ there. \qed
\end{proof}

\emph{Common input, restated.} At the good level, $|\nabla u|$ is close
to a constant on $\Sg=\Sg_{\hat t}$ \emph{in $L^2$} (a small absolute
constant, not a $\dT$-power). Nothing pointwise, nothing geometric, is
given for free. The two routes are two ways to use this.

\section{What Strategy~A must do, and the role of Part~A}

Strategy~A wants the \emph{geometric} conclusion that $\Sg$ is a small
graph over a sphere. Your proposed mechanism has two steps.

\emph{(A1) Self-quantification.} From the $L^2$ smallness of
$|\nabla u|-\fbar$ plus Markov, off a set of small $\Hn$-measure the
quantity $|\nabla u|$ is pinched into $[\fbar-\kappa,\fbar+\kappa]$.
Quantitatively, by Chebyshev on \eqref{eq:wv} (this is exactly
Lemma~\ref{lem:sg} below), the set $\{|\nabla u|\le s\}\cap\Sg$ has
measure $\lesssim s\,D_H/\fbar^2$.

\emph{(A2) Poisson smoothing.} Pass to an interior level
$\{u>t\}$, $t\ge C\dT$ with $C$ large, and use \emph{interior elliptic
regularity}: the value of a solution at interior distance $d$ from the
boundary is, by the Poisson representation, an average of boundary
data, so a boundary irregularity concentrated on a small set is
\emph{smoothed out} a definite distance inside. There is no dispersion;
this is the standard fact that
$\norm{\nabla^k(\text{harmonic ext})}_{L^\infty(\{\dist\ge d\})}
\lesssim d^{-k-(n-1)/p}\norm{\text{data}}_{L^p(\partial)}$.

The decisive question is the constant. If one estimates the interior
from a \emph{rough} $\partial\Om$, the factor $d^{-k-(n-1)/p}$ blows up
as $d\downarrow0$, and the good level sits in a $\dT$-thin shell, so
$d\sim\dT\to0$ and the estimate is vacuous. This is why the premise
matters: \textbf{Strategy~A presupposes Part~A}, namely that a prior
selection/regularization step has already produced a domain whose
boundary is smooth with \emph{$\dT$-uniform} $C^k$ bounds. \emph{Given
that premise}, global up-to-the-boundary Schauder makes
$u\in C^{k+1}(\overline\Om)$ with $\dT$-uniform norms; the interior
constant no longer depends on $d$, the small bad set is genuinely
smoothed on a slightly deeper level, and pointwise $|\nabla u|\approx
\fbar$ together with the controlled geometry bootstraps the higher
derivatives of the graph (an interpolation between small $L^\infty$ and
the fixed high norm). Under that premise the mechanism is sound; its
single hard point is the premise itself, i.e.\ whether the
selection/regularization genuinely delivers a $\dT$-uniform smooth
boundary. (That is precisely the unremoved compactness in the BDV
selection principle.)

\section{What Route~B does instead}

Route~B refuses the geometric conclusion. It keeps everything
\emph{interior} and in $L^2$, and it is Serrin's overdetermined problem
handled by Weinberger's argument, made quantitative.

\subsection{The Weinberger function}
For the torsion function $v$ of a domain $D$ (so $-\Delta v=1$, $v=0$
on $\partial D$), set $P:=|\nabla v|^2+\tfrac2n v$. Because
$\Delta v\equiv-1$ is constant,
\[
  \Delta P=\Delta|\nabla v|^2+\tfrac2n\Delta v
  =2|D^2v|^2+2\nabla v\!\cdot\!\nabla(\Delta v)-\tfrac2n
  =2\Bigl(|D^2v|^2-\tfrac1n\Bigr).
\]
The pointwise Cauchy--Schwarz inequality
$|D^2v|^2\ge(\tr D^2v)^2/n=(\Delta v)^2/n=1/n$ gives $\Delta P\ge0$;
equality at a point holds iff $D^2v=-\tfrac1n I$ there, i.e.\ iff $v$
agrees to second order with a ball quadratic $a-|x-x_0|^2/(2n)$.
Writing the trace-free Hessian one checks the exact identity
\begin{equation}\label{eq:DP}
  \Delta P=2\,\bigl|\,\TF\,\bigr|^2,\qquad
  \bigl|\TF\bigr|^2=|D^2v|^2-\tfrac1n\ge0,
\end{equation}
so the \emph{interior} quantity
$\R(D):=\int_D|\TF|^2$ measures non-roundness, vanishing iff $D$ is a
ball. Serrin/Weinberger rigidity is ``$\R(D)=0\Rightarrow D$ a ball'';
Route~B is its quantitative form, in two halves.

\subsection{Half 1: the deficit dominates $\R$}
We make the two classical identities explicit.

\emph{Rellich--Pohozaev.} Test $-\Delta v=1$ with $x\!\cdot\!\nabla v$.
Using $v=0$ (hence $\nabla v=v_\nu\nu$) on $\partial D$ and the two
divergence identities
\[
  \int_D\nabla v\!\cdot\!\nabla(x\!\cdot\!\nabla v)
  =\Bigl(1-\tfrac n2\Bigr)\!\int_D|\nabla v|^2
   +\tfrac12\!\int_{\partial D}|\nabla v|^2(x\!\cdot\!\nu),\qquad
  \int_D 1\cdot(x\!\cdot\!\nabla v)=-n\!\int_D v,
\]
together with $\int_D|\nabla v|^2=\int_D v=:T(D)$, one obtains after
cancellation
\begin{equation}\label{eq:poho}
  \tfrac12\int_{\partial D}|\nabla v|^2\,(x\!\cdot\!\nu)\,d\Hn
  \;=\;\frac{n+2}{2}\,T(D).
\end{equation}
Only the divergence theorem on a regular level is used --- no boundary
regularity of the original $\Om$.

\emph{Weinberger.} Integrate \eqref{eq:DP}. Pairing $\Delta P=2|\TF|^2$
against the Green function of $D$ (a nonnegative weight available with
no boundary regularity) converts the left side into $\int_D|\TF|^2$ up
to dimensional factors; the boundary terms produced by $\Delta P$,
using the boundary value $P=|\nabla v|^2$ on $\partial D$, are exactly
the deviation of $|\nabla v|$ from its Serrin constant. Combined with
\eqref{eq:poho} and the Saint--Venant gap
$T(B_D)-T(D)\ge0$ (a nonnegative boundary functional vanishing only on
balls), one gets:
\begin{proposition}\label{prop:half1}
For a regular level domain $D$ there is a purely dimensional
$c(n)>0$ with
\[
  \dT(D)\ \ge\ c(n)\,|D|^{-1}\,\R(D),
\]
no boundary regularity, no interior-ball or shape constant, additive
over connected components.
\end{proposition}
The pointwise Newton inequality $|D^2v|^2\ge1/n$, sharp only on the
ball quadratic, is exactly what makes the boundary and interior
expressions comparable and forces the relation to be \emph{order two}
(quadratic, the sharp Faber--Krahn exponent).

\subsection{Half 2: $\R$ controls the asymmetry}
\begin{lemma}[Hessian rigidity / quantitative Liouville]\label{lem:korn}
Let $D$ be open, $|D|=\omega_n$, $w\in H^2(D)$. There is a quadratic
$q(x)=a-|x-x_0|^2/(2n)$ and a constant $K(D)$ with
\[
  \int_D|\nabla w-\nabla q|^2\ \le\ K(D)\int_D\bigl|D^2w+\tfrac1n I\bigr|^2 .
\]
$K(D)$ is a Poincar\'e-type constant of $D$.
\end{lemma}
\begin{proposition}\label{prop:half2}
With $D$ a good level domain there is a dimensional $C(n)$ with
$\ \Asym(D)^2\le C(n)\,|D|^{-1}\,\R(D)$.
\end{proposition}
\begin{proof}[Mechanism]
$\R(D)$ small means $D^2v\approx-\tfrac1n I$ in $L^2$; by
Lemma~\ref{lem:korn}, $v$ is $H^2$-close to a ball quadratic $q$ on the
bulk; the superlevel set $\{v>\hat t\}$ is then close in measure to the
ball $\{q>\hat t\}$ at the sharp quadratic rate (an
$H^2\hookrightarrow$ trace estimate turns the interior closeness into
the squared $L^2$ size of the boundary deviation, which is
$\Asym^2$). The only delicate point is that $K(D)$ in
Lemma~\ref{lem:korn} could degrade where $\nabla v$ is small. This is
controlled by the small-gradient lemma below: that set has small
measure and is \emph{excised}, paid for additively.
\end{proof}

\subsection{The role of $D_H$: measure-control, not graphing}
\begin{lemma}[small-gradient set]\label{lem:sg}
On $\Sg=\Sg_{\hat t}$, for every $s>0$,
\[
  \Hn\bigl(\{x\in\Sg:|\nabla u(x)|\le s\}\bigr)
  \ \le\ \frac{4s}{\fbar^2}\int_{\Sg}\frac{(|\nabla u|-\fbar)^2}{|\nabla u|}\,d\Hn .
\]
\end{lemma}
\begin{proof}
Where $|\nabla u|\le s\le\fbar/2$ one has
$(|\nabla u|-\fbar)^2\ge\fbar^2/4$ and hence
$(|\nabla u|-\fbar)^2/|\nabla u|\ge\fbar^2/(4s)$; integrate
\eqref{eq:wv}'s integrand over that set only and divide. \qed
\end{proof}
By Lemma~\ref{lem:markov} the right side is a small absolute constant,
so the small-gradient set has small measure. \emph{This is the only
place $D_H$ is used in Route~B}: not to make $\Sg$ a graph, but to
measure-bound the set where the Hessian-rigidity constant could
degrade, so that it can be removed. For a disconnected good level one
argues per component: a component carrying order-one asymmetry has
order-one deficit by the qualitative Saint--Venant inequality (a
threshold, by compactness), hence lies outside the small-deficit
regime; so the per-component treatment loses nothing.

\subsection{Assembly}
Combining Propositions~\ref{prop:half2} and \ref{prop:half1} on the
good level domain, then the exact transfer $\dT(\text{level})\le
2\dT(\Om)$ (the shifted torsion function is exactly the level
domain's) and the elementary superlevel transfer
$\Asym(\Om)\le\Asym(\text{level})+2|\Om\setminus\text{level}|/|\Om|$
with the discarded layer $O(\dT)$, gives
$\Asym(\Om)^2\le C_*(n)\,\dT(\Om)$, the squared collar being
$O(\dT^2)\ll\dT$.

\section{The precise contrast}

\begin{itemize}
\item Strategy~A needs the bad set \emph{smoothed away}: pushed toward
measure zero and the good control made \emph{uniform and pointwise on
the whole boundary}, so that the level set is a controlled small graph.
That is the (A1)+(A2) program; its soundness rests entirely on Part~A
making the up-to-boundary Schauder constants $\dT$-uniform.
\item Route~B never needs the bad set to vanish or any pointwise
boundary control. It needs only that the bad set has \emph{small
measure} (true by Lemma~\ref{lem:sg}, with no propagation and no
Part~A), then \emph{deletes} it (Lemma~\ref{lem:korn} on the bulk) and
extracts the conclusion from the interior $L^2$ Hessian via
\eqref{eq:DP}. No graph, no boundary norm, no Poisson propagation.
\end{itemize}
So the smoothing-by-Poisson idea is the right tool \emph{for
Strategy~A}; Route~B's design point is that, arguing in the interior
through Weinberger, that step is simply not needed --- the same $L^2$
near-Serrin input is consumed by excision instead of by propagation.

\section{The Serrin dictionary}

Serrin's theorem: if the torsion function of $D$ has constant gradient
on $\partial D$, then $D$ is a ball. There are two quantitative
``stabilities'' of this statement, and the two routes are exactly the
two of them:
\begin{itemize}
\item \emph{Geometric/boundary Serrin stability} (Aftalion--Busca--Reichel,
Magnanini--Poggesi, Brandolini--Nitsch--Salani--Trombetti type):
$|\nabla v|$ near constant on $\partial D$ $\Rightarrow$ $\partial D$
close to a sphere \emph{as a graph}. This needs boundary regularity and
is what Strategy~A invokes.
\item \emph{Interior/Weinberger rigidity stability}: the integrated
trace-free Hessian $\R(D)$ is controlled by the deficit and controls
the asymmetry, with no boundary statement at all. This is Route~B.
\end{itemize}
Both descend from the same Weinberger $P$-function; they differ only in
whether one reads the equality case off the \emph{boundary} (needs
regularity) or off the \emph{interior} integral \eqref{eq:DP} (does
not). That is the entire conceptual content of ``Route~B uses Serrin
directly''.

\section{A worked illustration: ball plus thin tentacle}

Let $\Om=B_1\cup T$ with $T$ a tube of length $L=O(1)$ and cross-radius
$w\ll1$, $|T|\simeq\kappa_n w^{n-1}L$. On the tube the transverse
Dirichlet barrier $\phi(r)=(w^2-r^2)/(2(n-1))$ caps $u$:
$\sup_T u\simeq_n w^2$, independently of $L$. An energy expansion gives
\[
  \dT(\Om)\ \simeq_n\ |T|\ =\ \kappa_n w^{n-1}L,\qquad
  \Asym(\Om)\ \simeq_n\ |T|,
\]
and, since on the thin tube $D^2u\approx-\tfrac1{n-1}I_{n-1}$
transversally (not $\approx0$), $|D^2u+\tfrac1nI|^2=\Theta_n(1)$ per
unit volume there, so $\R(\Om)\simeq_n|T|$. Hence Route~B's chain reads
\[
  \Asym(\Om)^2\ \simeq\ |T|^2\ \ \ll\ \ |T|\ \simeq\ \R(\Om)\ \simeq\ \dT(\Om),
\]
so $\Asym^2\lesssim\dT$ holds \emph{with enormous room}: a tentacle is
\emph{not} a counterexample. By contrast the graph route must reach a
tentacle-free level; that costs discarding a ball-shell of measure
$\simeq_n w^2$ (the level $\{u>t\}$ for $t>\sup_T u\simeq w^2$ peels an
annulus of thickness $\simeq w^2$ off the ball part), while the budget
is $O(\dT)=O(w^{n-1}L)$. For $n\ge4$, $w^2\gg w^{n-1}L$: the graph
route cannot escape the tentacle within an $O(\dT)$ layer, whereas
Route~B never tries to --- it reads the conclusion off the interior
$\R$, where the tentacle is harmless. (Note: the tube-level
$D_H\simeq_n w^{n-3}L$ tends to $0$ for $n\ge4$, so such a tentacle
does \emph{satisfy} the near-Serrin hypothesis; this is exactly why
the graph route is genuinely obstructed there and the interior route
is not.)

\section{Honest status}

This note explains the \emph{design} --- why Route~B avoids the
Poisson-propagation step --- not a proof. Two statements are
load-bearing and not human-refereed:
\begin{checkpoint}[Half 1]
Proposition~\ref{prop:half1}: that the Saint--Venant deficit dominates
$|D|^{-1}\R(D)$ with a \emph{purely dimensional} constant, valid on a
regular level domain with no boundary regularity and stable under
disconnection and thin tentacles. The exact tracking of the
Green-function pairing constant is what must be checked.
\end{checkpoint}
\begin{checkpoint}[Half 2]
That the Poincar\'e constant $K(D)$ of Lemma~\ref{lem:korn} stays
\emph{dimensional} on the actual good level domain --- possibly
pinched, non-convex, multi-component --- after excising the
small-gradient set of Lemma~\ref{lem:sg}. This is the single most
important line to referee.
\end{checkpoint}
Symmetrically, Strategy~A's single hard point is the Part~A premise
(a $\dT$-uniform smooth boundary), which is the unremoved compactness
of the BDV selection principle. Neither route is refereed; they are
pursued in parallel because each has exactly one distinct analytic
crux of comparable standing, and the apparent internal coherence of
either is not, by itself, evidence --- only external scrutiny is.
\end{document}
